% Background rendering depends on graphicx and TikZ, which are loaded by
% layout/preamble.tex.  Its image and default-opacity settings live in
% config/style.tex, and are expanded when a page is shipped out.
\makeatletter
\newcommand{\MMBackgroundOpacity}{\DefaultBackgroundOpacity}
\newif\ifMMBackgroundOpacityOverridden
\MMBackgroundOpacityOverriddenfalse

% Render the configured image behind the full physical page.  The empty image
% path used by a style configuration disables backgrounds without an error.
\newcommand{\MMRenderPageBackground}{%
  \ifx\BackgroundImagePath\@empty
  \else
    \begin{tikzpicture}[remember picture, overlay]
      \node[opacity=\MMBackgroundOpacity] at (current page.center) {%
        \includegraphics[width=\paperwidth, height=\paperheight]{\BackgroundImagePath}%
      };%
    \end{tikzpicture}%
  \fi
}

% Use \changebackground{<opacity>} before the page break for the one page
% being completed.  Once that page has shipped out, the default opacity is
% restored automatically.
\newcommand{\changebackground}[1]{%
  \gdef\MMBackgroundOpacity{#1}%
  \global\MMBackgroundOpacityOverriddentrue
}

\AddToHook{shipout/background}{%
  \MMRenderPageBackground
  \ifMMBackgroundOpacityOverridden
    \global\def\MMBackgroundOpacity{\DefaultBackgroundOpacity}%
    \global\MMBackgroundOpacityOverriddenfalse
  \fi
}
\makeatother
